\section{Lessons and Threats to Validity}\label{sec:threats}

% ICDE: summary of findings and recommendations
\begin{table*}[t]
\centering\footnotesize
\caption{Summary of findings and recommendations.}\label{tab:findings}
\begin{tabularx}{\textwidth}{@{}l >{\raggedright\arraybackslash}X >{\raggedright\arraybackslash}X l@{}}
\toprule
& Finding (evidence) & Recommendation & Audience\\
\midrule
F1 & Under one SQLite certifier, released counterexamples lose about half of their pairs at the DDL tier and 82--92\,\% at the format tier; the ranking of tools changes with the tier (Sec.~\ref{sec:main}). & Compare disprovers only under one certifier, tier and budget; re-certify released witnesses on the benchmark's engine. & tool authors\\
F2 & By an LLM attribution, \V{attrImplMin}--\V{attrImplMax}\,\% of every tool's schema-tier witnesses rely on implausible data; humans shown facts of the data reach a similar verdict (\V{hvTwoImpConsOne}/\V{hvTwoConsOne} agreed items) but agree with the AI only moderately per item; humans without the facts almost never do; a repair test removes \V{xVanPct}\,\% of the differences (Sec.~\ref{sec:attrresults}). & Report witnesses with an attributed sample and the data facts the judgement rests on, never bare ``non-equivalent'' counts. & tool authors, maintainers\\
F3 & A data profile of the benchmark (minutes to mine) cuts implausible witnesses from \V{pCGImpl}\,\% to \V{pCGRImpl}\,\%; SpotIt+'s counterexamples (LLM-validated column constraints) are \V{pSIPImpl}\,\% implausible; re-certified counts are insensitive to the profile's thresholds (Sec.~\ref{sec:rresults}). & Publish the data profile with the benchmark; judge witnesses at $\mathcal A_R$, with a presupposition constraint for entities named in the question. & maintainers\\
F4 & The effect of witnesses on EX lies between \V{ptsErrRMin}--\V{ptsErrOnefMax} points per system (weighted by attributed error rates; a lower bound) and 15--20 points (all schema-tier witnesses counted) (Sec.~\ref{sec:rankings}). & Never subtract witness counts from EX; use meaningful witnesses to repair gold queries and questions. & maintainers, system builders\\
F5 & Join-graph completion and real-row retrieval matter; three synthetic components do not; a data-grounded audit of one system's accepted test predictions costs under one core-hour at 5\,s per pair, plus attribution (Sec.~\ref{sec:ablation}). & Build disprovers on the join graph and real data; certify and attribute every leaderboard submission. & tool authors, maintainers\\
F6 & On Spider, schema-tier witnesses again fail the profile, but CaliberGraph-R's lead does not replicate and fuzzing-R's witnesses are more meaningful; on 13 unseen BIRD systems the lead in counts replicates (Sec.~\ref{sec:replication}). & Mine profiles from instances large enough to expose keys and dependencies, or add them by hand. & maintainers, tool authors\\
\bottomrule
\end{tabularx}
\end{table*}

\textbf{Lessons.} Table~\ref{tab:findings} maps each finding to a recommendation. The common thread is that a counterexample is a database and should be judged as data: certified on the benchmark's engine with the original query strings, checked against the benchmark's own profile, and released so that others can attribute it. Meaningful witnesses are as much feedback on the benchmark as on the systems: 58--64\,\% of the sampled data-grounded witnesses expose latent gold-query defects or ambiguous questions (some on sub-databases without the named entity), against 11--22\,\% that expose prediction errors, and no correction of system scores can fix the former.

\textbf{Disproving is not proving.} A missing witness says nothing; accuracy corrected with witnesses is an upper bound on semantic accuracy, and error-weighted corrections are lower bounds on the over-statement. Bounded verifiers can also report the absence of counterexamples up to a bound, which our tools cannot.

\textbf{What ``implausible'' means.} The profile is mined from one instance, so it can be too strict (a plausible but unobserved value is admitted only through range and shape rules) and too weak: annotators still found conditional constraints it does not express (a badge awarded once while others repeat), which conditional-dependency and denial-constraint discovery could capture~\cite{fan2011cfd,chu2013dc}. ``Implausible'' combines absence from the instance with the annotators' judgement about the domain; a live football database may well contain unplayed fixtures without scores (Fig.~\ref{fig:example}). Sub-databases can also omit the entity a question names (Section~\ref{sec:rresults}). Some differences rely on behaviour SQL leaves unspecified, such as which tied row \texttt{LIMIT} keeps; BIRD's EX is computed with SQLite, so we count them, and attribution classifies most as ambiguities.

\textbf{AI-performed annotation.} All attribution and root-cause labels were produced by the AI model that also implemented the tools, in two independent passes plus adjudication; agreement between passes measures consistency, not correctness. A second model reproduces the labels (Section~\ref{sec:attrresults}), but both share a vendor and the same data access. The first author's earlier checks of two samples of 30 witnesses (agreement 26--28 of 30) were not blind: the AI verdicts were printed on the same sheet. The two human annotators, data analysts familiar with SQL who volunteered, labelled independently. The pre-registered blind check failed in its first round, in which they had no access to $D_0$; in a second, separately pre-registered round with facts from $D_0$ they agreed with each other and reproduced the difference between the tiers, but agreed with the AI only moderately item by item, so the AI's estimated shares carry this uncertainty. The two round-2 sheets returned to us contained a few manual edits outside the label fields (end-of-file markers, two altered witness cells in one item and a shortened fact line in another, both labelled alike by the two annotators; excluding the two items changes no outcome); the authors confirm that both rounds were annotated by hand without AI tools, and the sheets are released as received. Blinding to the tier is imperfect, since synthetic values (\texttt{'x1'}, \texttt{-1}) are recognisable. All annotations, reasons, witnesses and both human sheets are released.

\textbf{Baselines.} SpotIt and SpotIt+ are used through their released outputs because their constraint mining needs an LLM API unavailable in our environment; their budgets are larger than ours. ParSEval-R certifies nothing at $\mathcal A_R$ and is not a fair representative of symbolic generation there. The frozen first-stage fuzzing baseline accepted the first differing database without consulting the certifier, unlike every other tool; we re-ran exactly the affected pairs (132 distinct EX-correct, 259 EX-incorrect) with certifier-in-the-loop acceptance and the same seeds and report the corrected counts; its attribution sample was drawn before the correction.

\textbf{Deviations from the pre-registrations.} (i)~The attribution of first-stage witnesses was added after the first test run; it replaced a pre-registered residual audit of undisproved pairs. (ii)~The pre-registered manual spot check of rule-based root-cause labels was replaced by two AI passes after an AI check found the rule-based labels imprecise; no human check of root causes was made. (iii)~Before the second freeze, the definition of $\mathcal A_R$ in the second pre-registration was revised on development data: the observed-value condition was restricted to categorical columns, measure columns were excluded from unique column combinations, and the value-dependency class was added; the record of this revision was written after the second test run, but the changed files are covered by the freeze hashes. (iv)~The second stage was designed after seeing first-stage test attributions (Section~\ref{sec:protocol}). (v)~Stage~3 (Section~\ref{sec:replication}, Section~\ref{sec:attrresults}) was pre-registered before any of its data were seen; the repair test was added after the first human round (post hoc), its procedure was changed once after a 12-item smoke run, a build fix for two non-UTF-8 cells of one Spider database preceded all Spider runs, and part of one run shared the machine with analysis jobs. (vi)~The pre-registrations are timestamped by commits of our development repository, not by a public registry; the released snapshot lets anyone verify the frozen method and baseline code against the recorded hashes (six analysis-only scripts changed after the freezes are listed), and the commit history is available to the program committee. The complete list of amendments, with timestamps, is in the artifact repository.

\textbf{Generalisation.} All results concern SQLite, BIRD and Spider~1.0. The systems answer the same questions, which our intervals account for by resampling questions. The tier was conceived after first-stage test attributions had been seen; its only test on databases no design decision had seen, Spider, did not replicate CaliberGraph-R's lead, which replicated only on new systems over the BIRD databases. Spider also shows that a profile mined from small instances is less informative; larger enterprise schemas and other engines remain to be studied.
